\vfill\pagebreak
\section{Functions that change their arguments}
\label{sec:mutability:parameters}

A parameter holds a copy of its argument (Section~\ref{sec:functions:parameters}).
For a slice, that copy is a length and an address, and the parameter shares the
elements of the argument, as \token{seen} shared those of \token{primes}. A
function can therefore change the elements of a slice that its caller gives it,
if its parameter says so.

\lstinputlisting[style=coloredverbatim, caption={A function that changes the elements of a slice, \textit{examples/mutability/discount.yr}}, label=lst:mutability:discount]{examples/mutability/discount.yr}
\vspace{-10pt}%
\begin{lstlisting}[style=bashVerb]
[108, 72, 40]
\end{lstlisting}

The parameter \token{prices} of line~3 is declared \token{dmut}, as a variable
would be, and its type is \token{mut [mut i32]}: the function may change the
elements it receives. The call, on line~11, passes \token{alias prices}. The
parameter and the variable of \token{main} are two names for the same elements,
and \token{alias} marks the place where the second name is created, as in
Section~\ref{sec:mutability:alias}. Figure~\ref{fig:mutability:parameter} shows
the memory during the call.

\begin{figure}[h]
  \centering
  \begin{adjustbox}{max width=\linewidth}
    \begin{tikzpicture}
      \node[memcell] (mlen) at (0, 0) {3};
      \node[memaddr, right=0mm of mlen] (mptr) {};
      \node[cflab, left=3mm of mlen] (mn) {\token{prices}};
      \node[memframe, fit=(mn)(mptr)] (main) {};
      \node[memtitle] at ([xshift=2mm]main.north west) {main};

      \node[memcell] (dlen) at (0, -1.6) {3};
      \node[memaddr, right=0mm of dlen] (dptr) {};
      \node[cflab, left=3mm of dlen] (dn) {\token{prices}};
      \node[memcell] (pc) at (0, -2.6) {10};
      \node[cflab, left=3mm of pc] (pn) {\token{percent}};
      \node[memframe, fit=(dn)(dptr)(pn)] (discount) {};
      \node[memtitle] at ([xshift=2mm]discount.north west) {discount};

      \foreach \value [count=\i from 0] in {108, 72, 45} {
        \ifnum\i<2
          \node[memhit] (h\i) at (5.2 + \i * 1.0, -0.8) {\value};
        \else
          \node[memcell] (h\i) at (5.2 + \i * 1.0, -0.8) {\value};
        \fi
      }
      \draw[mempoint] (mptr.center) to[out=0, in=180] (h0.west);
      \draw[mempoint] (dptr.center) to[out=0, in=180] (h0.west);

      \node[memregion] (stack) at ($(main.north) + (0, 5mm)$) {stack};
      \node[memregion] at (h1 |- stack.base) {heap};
      \draw[memsep] (3.8, 1.0) -- (3.8, -3.3);
    \end{tikzpicture}
  \end{adjustbox}
  \caption{\label{fig:mutability:parameter} The call
    \token{discount(alias prices, 10)} of Listing~\ref{lst:mutability:discount},
    during the third turn of its loop. The call has a frame of its own, below
    that of \token{main}, and its parameter \token{prices} received the length
    and the address of the variable of \token{main}. The two designate the same
    elements: the first two are already changed.}
\end{figure}


The loop of line~4 goes through the indices of the slice, and line~5 replaces
each price with its value reduced by \token{percent} percent. The division is
between integers, so 45 reduced by 10\% gives 40, not 40.5. When the call
returns, its frame is given back, but the elements it changed are those of
\token{main}, which prints them.

The \token{alias} of the call is required. Without it, the call is refused.

\begin{lstlisting}[style=coloredverbatimError, escapechar=@, caption=A call that shares elements without \token{alias}]
use std::io;

fn discount(dmut prices: [i32], percent: i32) {
  for i in 0 .. prices.len {
    prices[i] = prices[i] * (100 - percent) / 100;
  }
}

fn main() {
  let dmut prices = copy [120, 80, 45];
  discount(@\hb{prices}@, 10); // error
  println(prices);
}
\end{lstlisting}

The compiler says \errmsg{E4226}{the call operator is not defined for
  \errhole{} and \errhole{}}, since no function \token{discount} accepts the
arguments as they are written, and the note under it gives the reason,
\textit{implicit alias of type mut [mut i32] is prohibited}. A reader of
\token{main} sees from the call alone that \token{discount} may change
\token{prices}: the \token{alias} is there. A call without it, such as
\token{average(prices)}, gives the function no right to change the elements.

\subsection{The parameter itself is still a constant}

\token{dmut} lets the function change the elements, not the parameter: giving
\token{prices} a new slice inside \token{discount}, \token{prices = copy [0];},
is refused with \errmsg{E4156}{prices is a value parameter, and cannot be used
  as a lvalue}. The new slice would only replace the copy of the length and the
address held by the frame of \token{discount}, and would have no effect on the
caller. Section~\ref{sec:mutability:references} shows how a function gives a new
value to a variable of its caller.

For the same reason, \token{mut} and \token{dmut} are refused on a parameter
whose value borrows nothing, such as \token{dmut n: i32} or
\token{dmut days: [i32 ; 12]}. Its argument is copied whole into the frame of
the function, so changing it could never affect the caller, and the compiler
says \errmsg{E4138}{a parameter cannot be mutable, if it is not a reference or
  does not borrow mutable data}.

\subsection{Returning elements that can change}

A function that creates a slice can return it with elements that can change,
with a return type written \token{dmut [i32]}.

\begin{lstlisting}[style=coloredverbatim]
fn zeros(n: usize)-> dmut [i32] {
  copy [0 ; n]
}

fn main() {
  let dmut counts = zeros(3);
  counts[1] = 5;
  println(counts);
}
\end{lstlisting}
\vspace{-10pt}%
\begin{lstlisting}[style=bashVerb]
[0, 5, 0]
\end{lstlisting}

The result of \token{copy} shares its elements with nothing, so it needs no
\token{alias}. Neither does the call, on the side of \token{main}: the function
has returned, and only \token{counts} designates the elements. A function that
returns the elements of one of its variables declared \token{dmut}, however,
returns them with \token{alias}, as in \token{alias result}: the variable and
the value returned designate the same elements, however briefly.
